\documentclass[a4paper,12pt]{book}
\usepackage[utf8]{inputenc}
\usepackage{graphicx}
\usepackage{amsmath}
\usepackage{amssymb}
\usepackage{cite}
\usepackage{enumitem}
\usepackage{mathptmx} 
\usepackage{titlepic}



\begin{document}

\begin{titlepage}
\titlepic{\includegraphics[width=\textwidth]{GoilGen_Logo.png}}
\author{Philipp Amrein}
\title{CoilGen Documentation}
\date{February 2022}
\end{titlepage}

%\frontmatter

\maketitle
\tableofcontents

\newpage

\section{{Purpose, Introduction}}
This \textit{CoilGen} project is supposed to be a community based tool that facilitates the design of coils within the MRI/NMR environment. A stream function approach of the Boundary Element Method is used to generate a current density to fit a freely specified target field. Further post processing is implement that processes the current density into a suitable coil layout.
Up to now, the code is written in MATLAB, but future migration to python might be advantageous, especially since it does not need proprietary software licenses. The author is very willing to collaborate with anyone who wants do the translation.

\section{Installation}
The project requires MATLAB and additionally FastHenry2 for calculation of the inductance.  The MATLAB version should not be older than \textit{2020A}. 

\section{Code description}
The algorithm is written within the main function named "coilgen.m". It can be called with various necessary and optional input parameters.



\subsection{File structure}

\begin{itemize}
	\item Documentation
	\item  Examples: Here are the Matlab scripts of the different projects that calls the main algorithm \textit{CoilGen.m} together with the input parameters
	\item  Geometry\textunderscore Data: \textit{.stl} files of surface geometries for the definition of the current carrying surfaces and the target fields. It should be in the bindary format.
	\item  sub\textunderscore functions: Folder of \textit{MATLAB} sub-functions that are called by the main function \textit{CoilGen.m} algorithm
	\item  Pre\textunderscore Optimized\textunderscore Solutions: \textit{.mat} data containers pre-optimized stream functions. Please check the function \textit{load\textunderscore preoptimized\textunderscore data.m}
	\item  plotting: Several functions to plot results.
	\item  Results: Folder for saving results.
	
\end{itemize}

\section{Usage}
"Coilgen.m" must be called with the required set of input parameters. All other settings and variables are optional. If not specified the default values set. The definitions of the default values can be seen within the function \textit{parse\textunderscore inputs.m}.

\subsection{Necessary input parameters}

\begin{itemize}
\item  \textit{coil\textunderscore mesh\textunderscore file}: The name of the \textit{.stl} file for the current carrying surface.  It must be in the\textit{ Geometry\textunderscore Data} folder.
\item  \textit{field\textunderscore shape}: The functional that specifies the target field i.e. the \textit{z} component of the magnetic field i.e. $B_{z}$. In general, it can be any function of the Cartesian coordinates $x,y,z$. To generate a linear transverse gradient specify: $field\textunderscore shape='y'$. This variable must be a character array. 

Other examples: 
\begin{itemize}
	\item  field\textunderscore shape='$x$' (linear x-Gradient)
	\item  field\textunderscore shape='$y$' (linear y-Gradient)
	\item  field\textunderscore shape='$z$' (linear z-Gradient)
	\item  field\textunderscore shape='$x.*sin(\alpha)-y.*cos(\alpha)$' (linear x-Gradient rotated with angle $\alpha$ around the z-axis)
	\item  field\textunderscore shape='$2.*x.*y$' (Spherical harmonic $S2$ shim field)
	\item  field\textunderscore shape='$1$' (constant field,  $B0$ offset)
\end{itemize}

The $x,y,z$ coordinates refer to the coordinate system which is also used for the coordinates of the surface geometry.
\textbf{Important}: Use the MATLAB $.*$ operator for element wise multiplication.

\end{itemize}

\newpage

\subsection{Optional input parameters}

All other settings and variables are optional. If not specified the default values set. The definitions of the default values can be seen within the function \textit{parse\textunderscore inputs.m}.


\begin{itemize}
	\item  \textit{pot\textunderscore offset\textunderscore factor}: Offset factor for contour levels, default value: 0,5, Type: numeric
	\item  \textit{target\textunderscore mesh\textunderscore file}: File of the target surface mesh, which will define the target field, default value: 'none', Type: Character array
	\item  \textit{ secondary\textunderscore target\textunderscore mesh\textunderscore file}:  File of a secondary target surface mesh, which will define the target field, default value: 'none', Type: Character array
	\item  \textit{secondary\textunderscore target\textunderscore weight}: weight for the secondary target points , default value: 1, Type: numeric
	\item  \textit{use\textunderscore only\textunderscore target\textunderscore mesh\textunderscore verts}: flag to use only the target mesh vertices as target coordinates, default value: false, Type: logical
	\item  \textit{sf\textunderscore source\textunderscore file}: file of an already optimized stream function, default value: 'none', Type: Character array
	\item  \textit{levels}: Number of potential levels, default value: 10, Type: numeric
	\item  \textit{level\textunderscore set\textunderscore method}: Specify one of the three ways the level sets are calculated:
	"primary","combined", or "independent", Default value: "primary" 0,5, Type: Character array
	\item  \textit{fieldtype\textunderscore to\textunderscore evaluate}: Fieldtype to evaluate; 'gradient' or 'field', default value: 'field', Type: Character array
	\item  \textit{surface\textunderscore is\textunderscore cylinder\textunderscore flag}: Fag for cylindrical surface; in case of cylinder, a special parameterization will be used: Default value: true, Type: Logical
	\item  \textit{circular\textunderscore diameter\textunderscore factor\textunderscore cylinder\textunderscore parameterization}: for the cylinder parameterization the ration of outer and inner boundary, default value: 1, Type: numeric
	\item  \textit{interconnection\textunderscore cut\textunderscore width}: The width in meter of the opening cut for the interconnection of the loops: Default value: 0.01, Type: numeric
	\item  \textit{target\textunderscore region\textunderscore radius}: Radius of a spherical target field, default value: 0.15, Type: numeric
	\item  \textit{target\textunderscore region\textunderscore resolution}: Number of target points per dimension within the target region: Default value: 10, Type: numeric
	\item  \textit{normal\textunderscore shift\textunderscore length}:  The distance in meter for which crossing lines will be seperated along the normal direction of the surface, Default value: 0.001, Type: numeric
	\item  \textit{min\textunderscore point\textunderscore loop\textunderscore number}: The minimal required number of point of a single loop; otherwise loops will be removed.., default value: 20, Type: numeric
	\item  \textit{area\textunderscore perimeter\textunderscore deletion\textunderscore ratio}: Additional loop removal criteria which relates to the perimeter to surface ratio of the loop, default value: 5, Type: numeric
	\item  \textit{max\textunderscore allowed\textunderscore angle\textunderscore within\textunderscore coil\textunderscore track}: Max allowed angle of the track of the contours, Default value: 120, Type: numeric
	\item  \textit{min\textunderscore allowed\textunderscore angle\textunderscore within\textunderscore coil\textunderscore track}: Min allowed angle of the track of the contours; smaller angles will be converted to straight lines in order to reduce the number of points, Default value: 0.0001, Type: numeric
	\item  \textit{tiny\textunderscore segment\textunderscore length\textunderscore percentage}: Minimum relative percentage for which points will be deleted which contribute to segments which is extremly short, default value: 0, Type: numeric
	\item  \textit{iteration\textunderscore num\textunderscore stream\textunderscore func\textunderscore refinement}: Number of refinement iterations of the mesh (together with the stream function), default value: 0, Type: numeric
	\item  \textit{b\textunderscore 0\textunderscore direction}: The direction (vector) along the interconnections will be aligned, default value: [0 0 1], Type: numeric array
	\item  \textit{track\textunderscore width\textunderscore factor}: In case of pcb layout, specify the track width: 0,5, Type: numeric
	\item  \textit{conductor\textunderscore cross\textunderscore section\textunderscore width}: cross section width of the conductor (for the inductance calculation) in meter, default value: 0.002, Type: numeric
	\item  \textit{conductor\textunderscore cross\textunderscore section\textunderscore height}: cross section height of the conductor (for the inductance calculation) in meter, default value: 0.002, Type: numeric
	\item  \textit{specific\textunderscore conductivity\textunderscore conductor}: Conducter conductiviy, default value:$  0.018*10^{-6}  $(Copper), Type: numeric
	\item  \textit{conductor\textunderscore thickness}: Thickness of the sheet current density of within the stream function representation: 0.005, Type: numeric
	\item  \textit{cross\textunderscore sectional\textunderscore points}: 2D edge points for direct defintion of the cross section of the conductor, default value: $ [2 1.5 ;2 -1.5; -2 -1.5; -2 1.5; 2 1.5 ]'/2000 $, Type: numeric array
	\item  \textit{geometry\textunderscore source\textunderscore path}: Directory of the .stl  geometry files: 'Geometry\textunderscore Data', Type: Character array
	\item  \textit{output\textunderscore directory}: Output directory for saving results, default value: 'Results', Type: character array
	\item  \textit{save\textunderscore stl\textunderscore flag}: Flag to save sweeped .stl: false, Type: logical
	\item  \textit{plot\textunderscore flag}: Flag to plot results, default value: true, Type: logical
	\item  \textit{interconnection\textunderscore method}: Interconnection method: Regular or spiral in/out, default value: 'regular', Type: Character array
	\item  \textit{skip\textunderscore postprocessing}: Flag to skip post processin, default value: false, Type: logical
	\item  \textit{skip\textunderscore inductance\textunderscore calculation}: Flag to skip inductance calculation, default value: false, Type: logical
	\item  \textit{force\textunderscore cut\textunderscore selection}: Force selection of the orientation of the interconnection areas \{'low' 'high'...\}, default value: \{\}, Type: cell
	\item  \textit{gauss\textunderscore order}: Gaus integration order, for the calculation of the sensitivity matrix default value: 2, Type: numeric
	\item  \textit{set\textunderscore roi\textunderscore into\textunderscore mesh\textunderscore center}: Flag to set the roi into the geometric center of the mesh, default value:false, Type: logical
	\item  \textit{tikonov\textunderscore reg\textunderscore factor}: Tikonov regularization factor for the SF optimization, default value: 1, Type: numeric
														
\end{itemize}

\newpage

\section{How to design your own coil with CoilGen}

Make a copy of \textit{y\textunderscore gradient\textunderscore coil.m} within the folder \textit{Examples} as a starting point. 

Wihtin this script, adjust the input parameters of the call of CoilGen to your need.  Especially, make sure the following parameters are well specified:

\begin{itemize}
	\item field\textunderscore shape: Choose the kind of target field: '$x$' ,'$y$' ,'$z$'  etc.
	\item coil\textunderscore mesh\textunderscore file: Specify the .stl file of the current carrying surface
	\item levels: Specifiy the number of turns of the coil
	\item surface\textunderscore is\textunderscore cylinder\textunderscore flag: This relates to the surface parameterization. If your current carrying is similar to a cylinder, set this flag to true
	\item interconnection\textunderscore cut\textunderscore width: The width of the connections, which determines their steepness
	\item normal\textunderscore shift\textunderscore length: The distance in meter for which crossing lines will be seperated along the normal direction of the surface
	\item iteration\textunderscore num\textunderscore stream\textunderscore func\textunderscore refinement: Choose a number between (0,1,2). The higher the number, the smoother the result, but the greater the computational effort.
\end{itemize}




\end{document}